% ============================================================================
% Portada y Resumen (Abstract)
% ============================================================================
\begin{titlepage}
  \thispagestyle{empty}
  \centering

  % Banda superior institucional
  {\color{unalverde}\rule{\textwidth}{2.5pt}}\\[0.35cm]
  {\large\bfseries\color{unalverde} UNIVERSIDAD NACIONAL DE COLOMBIA}\\[0.15cm]
  {\normalsize\color{gristexto} Facultad de Ingeniería, Curso de Inteligencia Artificial}\\[0.25cm]
  {\color{grisborde}\rule{\textwidth}{0.8pt}}

  \vspace{2.2cm}

  {\Large\color{azulclaro}\bfseries INFORME TÉCNICO DEL PROYECTO}\\[0.5cm]

  {\Huge\bfseries\color{azuldoc} Agente autónomo solucionador\\[0.25cm] de laberintos 2D en ESP32-S3}\\[0.6cm]

  {\large\itshape Hardware, control embebido con sensores IR y giroscopio,\\
  planificación a bordo y comparación de algoritmos de búsqueda}

  \vspace{1.6cm}

  \begin{tcolorbox}[
    colback=azulfondo,
    colframe=azuldoc,
    arc=4pt,
    boxrule=0.8pt,
    width=0.92\textwidth,
    halign=center
  ]
    \textbf{\color{azuldoc} Repositorio oficial del proyecto:}\\[2pt]
    \url{https://github.com/dabocanegrac/agente-laberinto}
  \end{tcolorbox}

  \vspace{1.4cm}

  {\normalsize\bfseries\color{azuldoc} INTEGRANTES DEL EQUIPO}\\[0.45cm]
  \begin{tabular}{ll}
    \toprule
    \textbf{Nombre completo} & \textbf{Correo electrónico} \\
    \midrule
    Daniel Alejandro Bocanegra Cifuentes & \texttt{dbocanegrac@unal.edu.co} \\
    Angel David Piñeros Sierra           & \texttt{apineross@unal.edu.co} \\
    Sara Sofía Rivera Sanabria           & \texttt{sariveras@unal.edu.co} \\
    Vanessa Valeria Robles Mendoza       & \texttt{varoblesm@unal.edu.co} \\
    Daniel Alonso Chavarro Chipatecua    & \texttt{dachavarro@unal.edu.co} \\
    \bottomrule
  \end{tabular}

  \vfill

  {\color{grisborde}\rule{0.7\textwidth}{0.6pt}}\\[0.3cm]
  {\normalsize\color{gristexto} Colombia, noviembre de 2026}\\[0.1cm]
  {\footnotesize\color{gristexto} Versión en curso: incluye pruebas ya corridas en pista y espacios marcados \emph{[POR COMPLETAR]} para los ensayos finales.}
\end{titlepage}

% Página de Resumen (Abstract)
\clearpage
\thispagestyle{empty}
\section*{Resumen}
\addcontentsline{toc}{section}{Resumen}

En este informe presentamos el diseño, la construcción y las pruebas de un agente físico que resuelve laberintos 2D, desarrollado para el curso de Inteligencia Artificial de la Universidad Nacional de Colombia~\cite{repo2026agentelaberinto,excel2026organizacion}. La meta era sacar los algoritmos clásicos de búsqueda del papel y correrlos en un carro diferencial de dos ruedas (2WD) con un microcontrolador ESP32-S3 N16R8, sin depender de un computador durante el recorrido.

Desde el principio descartamos un esquema teleoperado en el que el PC dicta cada movimiento (\texttt{DEC-05}). En su lugar, el PC solo envía el mapa en una grilla $(2n+1)\times(2m+1)$ de ceros y unos por Wi-Fi (TCP puerto~23) o por cable serial. A partir de ahí el carro trabaja solo: valida el archivo, lo guarda en memoria no volátil (\texttt{Preferences}/NVS), planea la ruta sobre el espacio de estados con orientación $\estado{f}{c}{o}$ y la ejecuta aunque se caiga el Wi-Fi. Como los motorreductores N20-10 no tienen \emph{encoders}, avanzar solo por tiempo acumula error rápido. Para corregirlo usamos tres sensores infrarrojos TCRT5000 (dos laterales para no salirse del carril y uno frontal para cuadrarse contra las paredes) junto con un giroscopio MPU6050 leído a \SI{500}{\hertz} en una tarea aparte de FreeRTOS, que controla los giros de \SI{90}{\degree} y mantiene el rumbo recto.

En el computador dejamos las herramientas de apoyo en Python (\texttt{pc/}): algoritmos de referencia, el generador de mapas aleatorios (\texttt{pc/generar.py}), el script que produce los planos a escala 1:1 y listas de corte (\texttt{pc/imprimir.py}), el comparador de métricas (\texttt{pc/comparar.py}) y una interfaz web para la demo (\texttt{pc/demo.py}). Sobre el mismo modelo de estados y costos ($c_F = 1.0$ al avanzar y $c_{\text{giro}}$ al girar) comparamos BFS, UCS, DFS, Greedy Best-First y A* con una heurística Manhattan ajustada por orientación que demostramos admisible y consistente. Además, compilamos el código C++ del carro dentro de \texttt{pytest} para comprobar que BFS (ya en \texttt{main}) y A* (en rama de integración) produzcan exactamente la misma ruta, costo y nodos expandidos que Python.

Recogemos aquí las decisiones tomadas durante el proyecto (\texttt{DEC-01} a \texttt{DEC-06}) y los cambios que surgieron al chocar con el hardware real, como el paso a baterías 2S~18650 con compensación por voltaje, el lazo PI para los giros y el comando \texttt{VOLVER}. También reportamos las pruebas preliminares ya verificadas en pista (6 de 6 corridas hasta la meta en el mapa de prueba de $4\times 4$, en unos \SI{14}{\second}) y dejamos la Sección~6 preparada con bloques \emph{[POR COMPLETAR]} para añadir los datos de los mapas físicos definitivos y los ensayos en exterior.

\vspace{0.6cm}
\noindent\textbf{Palabras clave:} Inteligencia Artificial aplicada, agente autónomo, planificación de rutas, búsqueda en espacio de estados, BFS, UCS, A*, ESP32-S3, MPU6050, TCRT5000, robótica móvil 2WD.
